\section{Who Uses Perricheno and How}
\label{sec:use_cases}

\sys{} is designed for any organisation or individual that needs to produce academic
or professional documents regularly, at high volume, or under tight time constraints.
Below are the four primary customer profiles we have identified from the current
user base.

\subsection{Universities and Research Institutes}

\subsubsection*{The situation}

Research universities in Kazakhstan and across the CIS are under increasing pressure
to publish.
Funding allocation, international rankings, and faculty promotions are all
partially or fully tied to publication counts and citation metrics.
At the same time, the administrative burden on researchers has grown: more grant
reports, more institutional documentation, more evaluation forms — all of which
require writing time that competes with research time.

\subsubsection*{How \sys{} helps}

\begin{itemize}
  \item \textbf{Conference paper drafts.}
        A researcher submitting to a conference can use \sys{} to produce the
        initial draft (introduction, related work, methodology outline) in under
        an hour, then focus their time on the results and discussion sections
        that require their specific expertise.
  \item \textbf{Grant application supporting documents.}
        Funding applications require extensive background sections that are
        labour-intensive to write but largely formulaic.
        \sys{} can produce these sections with a verified bibliography in
        a fraction of the time.
  \item \textbf{Student thesis assistance.}
        Graduate students can use \sys{} to generate a reference-verified
        literature review chapter, which is typically the most time-consuming
        chapter to write and the one most vulnerable to citation errors.
  \item \textbf{Institutional reports.}
        Annual research output reports, accreditation documentation, and
        departmental reviews can be drafted at scale.
\end{itemize}

\subsubsection*{Typical productivity gain}

Based on feedback from the Astana IT University pilot:
researchers using \sys{} for literature review and background sections reported
saving 2--4 hours per document.
Across a department of 20 active researchers producing 3--5 documents per year,
this translates to 120--400 researcher-hours saved annually.

\subsection{Corporate R\&D and Innovation Departments}

\subsubsection*{The situation}

Corporate research teams — in technology, pharmaceuticals, finance, energy —
produce technical reports, white papers, feasibility studies, and patent backgrounds
that share the same structure as academic papers: referenced claims, methodologies,
and results presented in a professional format.
These teams typically do not have dedicated technical writers; scientists and
engineers write their own documentation.

\subsubsection*{How \sys{} helps}

\begin{itemize}
  \item \textbf{Technology landscape reports.}
        A team exploring a new technology area can use \sys{} to produce a
        survey of the current state of the art — with verified academic sources —
        as a starting point for internal discussion.
  \item \textbf{Patent prior art documentation.}
        Patent applications require a background section demonstrating knowledge
        of prior art.
        \sys{} can generate a reference-verified prior art survey from a topic
        specification.
  \item \textbf{Competitive analysis briefings.}
        Technical briefings for management that require academic citations
        to support claims about technology readiness or market trends.
  \item \textbf{Regulatory submission support.}
        In regulated industries (pharmaceuticals, medical devices), submissions
        must include extensive literature reviews.
        \sys{} can produce the initial draft of these sections at scale.
\end{itemize}

\subsection{Academic Publishers and Editing Services}

\subsubsection*{The situation}

Academic publishers and editing companies receive manuscripts that frequently
have citation problems: broken references, incorrectly formatted citations,
missing DOIs, and outright fabricated sources.
Checking these manually is slow and expensive.

\subsubsection*{How \sys{} helps}

Publishers can use the \sys{} verification layer as a standalone reference
audit tool: submit a document's bibliography, and the platform checks every
entry against live academic databases, flags unresolvable entries, and
classifies each reference by authority level (peer-reviewed journal,
conference proceedings, preprint, web resource, or unresolvable).
This turns a multi-day manual audit into a minutes-long automated process.

\subsection{Individual Researchers and Graduate Students}

\subsubsection*{The situation}

An individual researcher working under publication pressure faces the same
problems as an institutional team, but without institutional resources.
They cannot afford a professional writing service.
They do not have time to manually verify 30 references.
They may not be fluent in \LaTeX{} or in the citation style required by their
target venue.

\subsubsection*{How \sys{} helps}

\sys{} lowers the barrier to professional academic document production.
A graduate student who has experimental results but struggles with the writing process
can use \sys{} to produce a structured first draft with a verified bibliography,
then revise it with their own voice and their specific contributions.
The result is a document that meets professional formatting standards from the start,
reducing the revision cycle with supervisors and reviewers.

\begin{keybox}[title=Real Session Example: Master's Thesis Literature Review]
\textbf{Topic:} ``Machine learning approaches to early detection of Type 2 diabetes''\\
\textbf{Language:} English\\
\textbf{Session time:} 34 minutes\\
\textbf{Output:}
\begin{itemize}[noitemsep]
  \item 12-page literature review with 28 verified references (all resolvable to live DOIs)
  \item 3 summary figures (comparison chart, timeline, method taxonomy diagram)
  \item IEEE-formatted bibliography compiled without errors
\end{itemize}
\textbf{Student's assessment:} ``This would have taken me a week to do manually.
The references were the biggest relief — I didn't have to check each one.''
\end{keybox}

\begin{figure}[h!]
\centering
\begin{tikzpicture}[
  ubox/.style={draw=#1, fill=#2, rounded corners=5pt,
               minimum width=5.5cm, minimum height=2.4cm,
               align=center, font=\scriptsize\bfseries, line width=1.4pt,
               text width=5.2cm, inner sep=8pt},
  hub/.style={draw=C1, fill=L1, circle, minimum size=2.4cm,
              align=center, font=\small\bfseries, line width=2pt},
  arr/.style={-{Stealth[length=6pt,width=4pt]}, line width=1.2pt, draw=DG!40},
]
\node[hub] (hub) at (0,0) {\sys{}\\Platform};

\node[ubox={C1}{L1}] (U1) at (0, 3.4)
  {Universities \& Research Institutes\\[0.3em]
   {\normalfont\tiny Conference drafts · Grant documents · Student theses}};

\node[ubox={C2}{L2}] (U2) at (5.8, 0)
  {Corporate R\&D\\[0.3em]
   {\normalfont\tiny Technology surveys · Patent prior art · Regulatory docs}};

\node[ubox={C3}{L3}] (U3) at (0,-3.4)
  {Academic Publishers\\[0.3em]
   {\normalfont\tiny Reference audit · Citation verification · Manuscript checking}};

\node[ubox={C5}{L5}] (U4) at (-5.8, 0)
  {Individual Researchers\\[0.3em]
   {\normalfont\tiny Literature reviews · Thesis chapters · Journal submissions}};

\draw[arr] (hub)--(U1);
\draw[arr] (hub)--(U2);
\draw[arr] (hub)--(U3);
\draw[arr] (hub)--(U4);
\end{tikzpicture}
\caption{Primary customer categories for the \sys{} platform.
         Each category uses the same underlying pipeline but for different
         document types and at different volumes.
         Universities and corporate R\&D teams are the highest-volume users;
         individual researchers represent the largest number of accounts.}
\label{fig:use_cases}
\end{figure}

